\documentclass[lettersize,journal]{IEEEtran}
\usepackage{amsmath,amsfonts}
\usepackage{algorithmic}
\usepackage{algorithm}
\usepackage{array}
\usepackage[caption=false,font=normalsize,labelfont=sf,textfont=sf]{subfig}
\usepackage{textcomp}
\usepackage{stfloats}
\usepackage{url}
\usepackage{verbatim}
\usepackage{graphicx}
\usepackage{booktabs}
\usepackage[noadjust]{cite}
\usepackage{hyperref}
\hypersetup{
    colorlinks=true,
    linkcolor=black,
    citecolor=blue,
    urlcolor=blue
}
\hyphenation{op-tical net-works semi-conduc-tor IEEE-Xplore}

\begin{document}

\title{SpatialGDC: a novel spatial domain identification method for spatially resolved transcriptomics data based on spatial-guided dual-graph contrastive learning}

\author{Haoze~Li,~\IEEEmembership{Student Member,~IEEE,}
        % TODO: Add affiliation and manuscript date before submission
        % \thanks{Manuscript received XXX; revised XXX.}
        % \thanks{H. Li is with the XXX, XXX University, City, Country.}
}

% The paper headers
\markboth{IEEE Transactions on Computational Biology and Bioinformatics}%
{Li \MakeLowercase{\textit{et al.}}: SpatialGDC: Spatial Domain Identification via Spatial-Guided Dual-Graph Contrastive Learning}

\maketitle

\begin{abstract}
Spatial transcriptomics provides a powerful perspective for resolving microenvironmental heterogeneity in complex tissues, where spatial domain identification serves as a critical foundational task. However, current graph neural network (GNN)-based methods face two key bottlenecks when processing spatial transcriptomics (ST) data: (1) feature graphs constructed from highly sparse expression matrices harbor long-range spurious connections, causing GNNs to erroneously aggregate distant interfering signals that pollute local biological features; and (2) standard contrastive learning strategies forcibly repel spatially adjacent homogeneous cells as negative samples (false negatives), disrupting the continuity of anatomical structures. To address these challenges, we propose \textbf{SpatialGDC}, a spatial domain identification framework based on spatial-guided dual-graph contrastive learning. SpatialGDC introduces a spatial-prior guided adaptive pruning mechanism that dynamically filters long-range noisy edges in the feature graph using physical distance as a prior probability, blocking the propagation of detrimental information. To prevent excessive repulsion among locally homogeneous cells, the model introduces a false-negative-aware contrastive loss, which imposes an exponential soft penalty on the shared neighborhoods of the dual graphs. Extensive benchmarks across 10x Visium, Stereo-seq, and Slide-seqV2 platforms demonstrate that SpatialGDC outperforms existing state-of-the-art methods, achieving a mean ARI of 0.59 on the DLPFC dataset and 0.63 on the BRCA dataset.
\end{abstract}

\begin{IEEEkeywords}
Spatial Transcriptomics; Spatial Domain Identification; Graph Contrastive Learning; Dual-view Modeling; Spatial-prior Guided Pruning; Graph Denoising.
\end{IEEEkeywords}

\section{Introduction}
\label{sec:introduction}

\IEEEPARstart{T}{he} recent advances in Spatial Transcriptomics (ST) have enabled researchers to profile gene expression patterns within tissues at high resolution while preserving the spatial topological context \cite{stahl2016visualization}. Compared to traditional single-cell RNA sequencing (scRNA-seq), which inevitably loses the spatial location information of cells during tissue dissociation \cite{stuart2019comprehensive}, ST technologies add a new dimension to genomics \cite{larsson2021spatially} by linking molecular profiles to their original anatomical positions. They provide valuable insights into spatial intercellular interactions \cite{rao2021exploring}, the heterogeneity of tissue microenvironments \cite{zhao2021cancer}, and the spatial dynamics of organ development \cite{asp2019spatiotemporal}. Among the various downstream analysis tasks of ST data, Spatial Domain Identification is the most fundamental and critical step for deconstructing complex tissue architectures \cite{wang2024comprehensive}.

However, accurate spatial domain identification still faces significant computational challenges. Due to the physical limitations of sequencing technologies, ST data ubiquitously exhibits extreme gene expression sparsity and the dropout phenomenon, resulting in severe technical noise \cite{qiu2020embracing}. Moreover, complex tissue structures often present an inconsistency between transcriptomic semantics and physical spatial topology \cite{dong2022deciphering}. Cells of the same type may be distributed across physically disjoint regions, whereas physically adjacent cells may belong to entirely distinct functional hierarchies. These two issues are coupled: noisy expression profiles make it harder to distinguish whether spatially proximate cells share genuine biological similarity or merely coincidental correlation. Overcoming expression sparsity while balancing these two modalities remains a core challenge \cite{zhao2021bayesspace}.

Existing mainstream spatial clustering methods exhibit obvious limitations in addressing the aforementioned challenges. Early traditional clustering methods (e.g., \textbf{Seurat} \cite{stuart2019comprehensive}) rely solely on the gene expression view, completely ignoring the spatial context. To introduce spatial constraints, subsequent methods have widely adopted Graph Neural Networks (GNNs) \cite{wu2020comprehensive}. The foundational graph architectures, such as graph convolutional networks \cite{kipf2016gcn}, aggregate local neighborhoods but face limitations. Certain models based on a single spatial graph, such as \textbf{SpaGCN} \cite{hu2021spagcn}, \textbf{STAGATE} \cite{dong2022stagate}, and \textbf{CCST} \cite{li2022cell}, over-rely on physical proximity, which tends to cause feature over-smoothing at tissue boundaries where different domains meet \cite{li2018deeper}. Multi-modal methods such as \textbf{Spatial-MGCN} \cite{wang2023spatialmgcn} and \textbf{stGRL} \cite{lu2025stgrl} have attempted to utilize dual-view modeling to bridge spatial and transcriptional domains. However, they often ignore the ``long-range false edge noise'' during the construction of the feature graph---edges that connect spatially distant spots with spuriously high expression correlation arising from technical noise rather than genuine biological similarity. These false edges inject noise into the GNN message-passing process, and more critically, the inherent assumption of the standard InfoNCE objective that all non-identical samples are equally negative causes the model to erroneously push apart adjacent homogeneous cells, disrupting the continuity of local structures \cite{chen2020measuring}.

To overcome the inherent defects of existing multi-modal methods, we argue that merely constructing a ``Dual-view'' architecture is far from sufficient. While standard Contrastive Learning paradigms can extract highly discriminative features by maximizing agreement between complementary views \cite{chen2020simple}, directly applying it to spatial graphs raises two problems that must be addressed jointly. First, the feature graph is inherently disrupted by long-range spurious edges due to expression noise, and a uniform random edge dropout treats all edges equally regardless of their biological plausibility. Without incorporating spatial distance as a prior, the model cannot distinguish between a spurious correlation across distant regions and a genuine biological connection between nearby spots. Second, the standard InfoNCE objective assumes that all non-identical samples are equally negative, an assumption that breaks down in spatially structured tissues where physically proximate cells frequently belong to the same domain. Repelling these false negatives as if they were truly dissimilar contradicts the spatial continuity that the model should preserve. Therefore, introducing physical spatial location as a prior to guide dynamic denoising of the feature view is crucial, and simultaneously designing a constraint mechanism that can perceive and tolerate tissue-homogeneous neighbors is the pathway to breaking current performance bottlenecks.

Motivated by these insights, this paper proposes a framework named \textbf{SpatialGDC} (Spatial-Guided Dual-graph Contrastive learning). We employ a spatial-guided dual-graph contrastive learning framework that disentangles physical spatial topology and transcriptional semantics through a shared-weight (siamese) GCN encoder with attention-based adaptive gating fusion. To address long-range noise in feature graphs, we design a spatial-prior guided adaptive pruning mechanism that uses a Gaussian decay function of physical distance to compute per-edge retention probabilities, performing eccentric edge dropout while preserving graph connectivity. To maintain tissue continuity, we introduce a false-negative-aware soft penalty contrastive loss that reduces the repulsion weight of intersection-neighborhood false negatives, avoiding mutual repulsion between homogeneous cells while still preserving their gradient signal for regularization. Extensive results demonstrate that SpatialGDC achieves a mean ARI of 0.59 on DLPFC and 0.63 on BRCA, outperforming existing methods on 10x Visium \cite{maynard2021dlpfc}, Stereo-seq \cite{chen2022stereoseq}, and Slide-seqV2 \cite{stickels2021slideseq} platforms.

\section{Related Work}
\label{sec:related_work}

\subsection{Spatial Domain Identification Methods}
Spatial domain identification aims to partition spots with similar transcriptomic profiles and spatial contexts into continuous functional regions, which is a fundamental task in spatial transcriptomics (ST) data analysis. Early computational methods predominantly directly adopted classic clustering algorithms from the single-cell RNA sequencing (scRNA-seq) era. Powerful toolkits such as \textbf{Seurat} \cite{stuart2019comprehensive}, \textbf{SCANPY} \cite{wolf2018scanpy}, and deep generative models like \textbf{scVI} \cite{lopez2018deep} were widely used. However, these methods solely utilize the gene expression matrix and completely ignore spatial topological information. To overcome this limitation, spatial statistical models and dimensionality reduction frameworks were introduced, including Hidden Markov Random Fields (HMRF) \cite{zhu2018hmrf}, \textbf{BayesSpace} \cite{zhao2021bayesspace}, \textbf{SpatialPCA} \cite{shang2022spatialpca}, and variance-based methods like \textbf{SPARK} \cite{sun2020spark}, which impose local spatial smoothness constraints from a probabilistic perspective. In recent years, deep learning methods have significantly improved the accuracy of spatial domain identification. Autoencoder-based approaches such as \textbf{SEDR} \cite{fu2021sedr} learn low-dimensional latent representations through reconstruction objectives, but their encoder architectures do not explicitly model spatial adjacency. For instance, \textbf{SpaGCN} \cite{hu2021spagcn} and \textbf{DeepST} \cite{xu2022deepst} construct spatial adjacency graphs based on physical coordinates to extract joint features. However, relying entirely on gene expression makes models highly susceptible to sequencing noise \cite{qiu2020embracing}, whereas relying solely on physical adjacency easily triggers feature over-smoothing at tissue boundaries \cite{li2018deeper}, making it difficult to accurately delineate complex anatomical edges.

\subsection{Graph-based Methods in Spatial Transcriptomics}
Graph Neural Networks (GNNs) \cite{wu2020comprehensive}, due to their inherent advantages in modeling non-Euclidean data topologies \cite{kipf2016gcn}, have become the mainstream framework for processing ST data. Regarding graph construction, single-graph methods such as \textbf{STAGATE} \cite{dong2022stagate} and \textbf{CCST} \cite{li2022cell} typically build a unitary spatial adjacency graph based on the physical coordinates of the spots to aggregate local microenvironment information. In contrast, multi-graph (or multi-view) methods attempt to better exploit multi-modal information \cite{zhang2023graphst}. For example, \textbf{Spatial-MGCN} \cite{wang2023spatialmgcn} and \textbf{stGRL} \cite{lu2025stgrl} attempt to simultaneously construct a feature expression graph and a spatial graph, conducting joint learning. \textbf{GraphST} \cite{zhang2023graphst} further integrates contrastive learning with spatial graph neural networks to learn discriminative spot representations. Although multi-graph modeling provides richer perspectives, existing methods exhibit obvious limitations. Specifically, feature graphs constructed based on expression correlation are extremely sensitive to sequencing depth and the dropout phenomenon \cite{jiang2022statistics}, easily generating long-range false connections. While theoretical GNN research has proposed techniques like \textbf{DropEdge} \cite{rong2019dropedge} to alleviate topological noise, existing ST multi-graph networks often lack dynamic, spatial-prior-guided denoising mechanisms. Blind information passing inevitably leads to distant technical noise polluting local biological features.

\subsection{Contrastive Learning in Bioinformatics}
Contrastive Learning, as a powerful unsupervised representation learning paradigm, has evolved rapidly from computer vision algorithms (e.g., \textbf{MoCo} \cite{he2020momentum}) to graph representation learning (e.g., \textbf{GCC} \cite{qiu2020gcc}). It extracts highly discriminative features by optimizing alignment and uniformity on the hypersphere \cite{oord2018infonce}. Recently, several methods have introduced graph contrastive learning into the ST spatial domain identification task \cite{you2020graphcl}. However, when directly migrating the standard InfoNCE contrastive loss to ST data, the construction of positive and negative samples faces severe domain mismatch challenges, since the assumption that all non-identical samples are equally negative does not hold in spatially structured tissues. Traditional hard negative mining strategies \cite{robinson2020contrastive} treat all other spots within a batch as negative samples and forcibly repel them. Nevertheless, in real biological tissues, physically adjacent cells often exhibit high spatial homogeneity. Blindly pushing away these high-confidence physical neighbors causes a severe ``False Negatives'' repulsion phenomenon \cite{chuang2020debiased}, thereby disrupting the inherent continuity of the tissue structure.

\subsection{Summary and Limitations of Existing Methods}
In summary, although existing mainstream algorithms have made significant progress in the spatial domain identification task, they have not yet thoroughly resolved the two core pain points of ST data. First, when processing multi-view graph structures, existing methods lack effective spatial prior interventions targeting long-range technical noise in the feature graph, making it difficult to achieve accurate pruning while maintaining structural connectivity. Second, regarding cross-view feature alignment, traditional graph contrastive learning strategies fail to consider the highly homogeneous microenvironment characteristics of spatial transcriptomes; the lack of tolerance for false negatives leads to impaired continuity of clustering boundaries. For instance, \textbf{SpatialDG} \cite{spatialdg2026} employs dual-graph Deep Graph Infomax with hard negative selection, yet its adaptive negative sampling may exclude biologically meaningful false negatives entirely rather than softening their repulsion. To address these limitations, this paper proposes the SpatialGDC framework to dynamically filter out long-range pseudo-connections and maintain the homogeneous microenvironment, achieving higher-precision spatial domain identification.



\section{Methods}
\label{sec:methods}

\subsection{Overview}
This study proposes SpatialGDC, a spatial transcriptomics analysis framework based on dual-graph contrastive learning with spatial-prior guided denoising. Given a spatial transcriptomics dataset containing $N$ spots, we extract its gene expression feature matrix $X \in \mathbb{R}^{N \times d_{in}}$ and spatial coordinate matrix $S \in \mathbb{R}^{N \times 2}$. As illustrated in Fig.~\ref{fig:architecture}, the core workflow of SpatialGDC comprises six modules: (A) \textbf{Data Inputs \& Graph Construction}: The spatial adjacency graph and the feature expression graph are constructed from spatial coordinates and gene expression profiles, respectively; (B) \textbf{Perturbation (Noise \& Pruning)}: Data augmentation is performed on both graphs and the input features through edge dropout and noise injection, where the feature graph employs spatially-informed non-uniform pruning guided by physical distance; (C) \textbf{Dual-Graph Encoders (GCNs)}: A shared-weight Graph Convolutional Network (GCN) encoder processes the two views independently---the same encoder parameters are applied to each graph topology, extracting spatial embeddings from $A_s'$ and expression embeddings from $\tilde{A}_f$ in two forward passes; (D) \textbf{Fusion \& Decoder (Reconstruction)}: An attention block adaptively fuses the dual-view embeddings, and a decoder reconstructs the input features via transposed GCN weight matrices; (E) \textbf{Multi-level Contrastive Learning (Multi-level CL)}: Instance-level and cluster-level contrastive learning are jointly employed to align cross-view representations, with a soft penalty mechanism to tolerate false negatives; (F) \textbf{Downstream Analysis Tasks}: The learned fused embeddings can be applied to a variety of downstream analytical tasks in spatial transcriptomics, including spatial domain identification, spatially variable gene detection, trajectory inference, and data denoising and imputation.

\begin{figure*}[htbp]
    \centering
    \includegraphics[width=\textwidth, trim=0cm 0cm 0cm 0cm, clip]{figures/fig_architecture.pdf}
    \caption{\textbf{Schematic diagram of the overall SpatialGDC architecture.} The framework comprises six modules: (A) Data Inputs \& Graph Construction from spatial coordinates and gene expression; (B) Perturbation (Noise \& Pruning) with uniform dropout on the spatial graph and Gaussian-decay pruning on the feature graph; (C) Dual-Graph Encoders (shared-weight GCN) processing spatial and expression graph topologies through the same encoder; (D) Fusion \& Decoder (Reconstruction) via adaptive attention weighted combination and transposed GCN decoder; (E) Multi-level Contrastive Learning (Multi-level CL) with instance-level soft penalty and cluster-level alignment; (F) Downstream Analysis Tasks including spatial domain identification, spatially variable gene detection, trajectory inference, and data imputation.}
    \label{fig:architecture}
\end{figure*}

\subsection{Data Preprocessing and Dual-Graph Construction}
Prior to model training, raw gene expression matrices undergo standardized preprocessing to reduce technical noise and computational burden. All preprocessing steps are implemented using the Scanpy package \cite{wolf2018scanpy}.

First, we filter out genes that are not expressed in any spot, obtaining the filtered count matrix $Y \in \mathbb{R}^{N \times g}$, where $g$ is the number of retained genes. Next, we perform library-size normalization to account for differences in sequencing depth across spots, followed by log-transformation to stabilize variance:
\begin{equation}
    \hat{Y}_{ij} = \log \left( \frac{10^4}{\sum_{k=1}^{g} Y_{ik}} \cdot Y_{ij} + 1 \right)
\end{equation}

To focus on genes with the strongest biological signal, we select the top 3,000 highly variable genes (HVGs) based on their standardized dispersion after variance-stabilizing transformation. Specifically, genes are binned by mean expression, and those with the highest dispersion within each bin are retained. This HVG selection step reduces dimensionality while preserving genes that carry the most informative expression variation across spots.

The HVG-filtered expression matrix is then standardized to zero mean and unit variance per gene:
\begin{equation}
    \tilde{Y}_{ij} = \frac{\hat{Y}_{ij} - \mu_j}{\sigma_j}
\end{equation}
where $\mu_j$ and $\sigma_j$ denote the mean and standard deviation of gene $j$ across all spots, respectively.

Finally, we apply Principal Component Analysis (PCA) to the standardized matrix and retain the top $d_{in}=200$ principal components as the input feature matrix $X \in \mathbb{R}^{N \times d_{in}}$ for downstream graph construction and model training. PCA serves a dual purpose: (i) it further reduces dimensionality and filters out residual technical noise; and (ii) it decorrelates the features, which stabilizes the subsequent calculation of Pearson correlation in the feature graph construction.

Spatial transcriptomics data inherently contains two modalities of information: the semantic similarity of gene expression and the spatial proximity of physical locations. Therefore, we construct a spatial view and a feature view, respectively.

For the \textbf{Spatial Graph}, we first compute the pairwise inverse Euclidean distance based on the spatial coordinate matrix $S \in \mathbb{R}^{N \times 2}$:
\begin{equation}
    w_{ij} = \frac{1}{\| s_i - s_j \|_2 + \epsilon}, \quad i, j = 1, 2, \dots, N
\end{equation}
where $\epsilon$ is a small constant for numerical stability. The spatial adjacency matrix $A_s \in \mathbb{R}^{N \times N}$ is then obtained by retaining the top-$K$ nearest neighbors per spot. The resulting weight matrix is symmetrized, binarized, and the diagonal entries are set to zero.

For the \textbf{Feature Graph}, we compute the Pearson correlation between spots based on the PCA-reduced expression matrix $X \in \mathbb{R}^{N \times d_{in}}$:
\begin{equation}
    \rho_{ij} = \frac{\sum_{m=1}^{d_{in}} (X_{im} - \bar{X}_i)(X_{jm} - \bar{X}_j)}{\sqrt{\sum_{m=1}^{d_{in}} (X_{im} - \bar{X}_i)^2} \sqrt{\sum_{m=1}^{d_{in}} (X_{jm} - \bar{X}_j)^2}}
\end{equation}
where $\bar{X}_i = \frac{1}{d_{in}} \sum_{m=1}^{d_{in}} X_{im}$. The feature adjacency matrix $A_f \in \mathbb{R}^{N \times N}$ is constructed by retaining the top-$k$ ($k=12$) most correlated neighbors per spot. The continuous correlation weights are symmetrized via $(W + W^\top)/2$, the diagonal is set to zero, and finally binarized with a correlation threshold of $0.1$ to filter out weakly correlated edges.

Following the standard GCN renormalization scheme~\cite{kipf2016gcn}, both adjacency matrices are first augmented with weighted self-loops and then symmetrically normalized:
\begin{equation}
    \hat{A} = \tilde{D}^{-1/2} \tilde{A} \tilde{D}^{-1/2}
\end{equation}
where $\tilde{A}$ is the self-loop-augmented adjacency and $\tilde{D}$ is its diagonal degree matrix. This renormalization prevents numerical instabilities and ensures balanced feature propagation across nodes with varying degrees.

\subsection{Perturbation (Noise \& Pruning)}
\label{sec:perturbation}
To improve model robustness and support contrastive learning, we design a three-pronged perturbation strategy (Module B: Perturbation in Fig.~\ref{fig:architecture}) that augments each input modality with targeted noise injection.

\textbf{Input Feature Perturbation.}
Before being fed into either encoder, the gene expression input $X$ is augmented with additive Gaussian noise followed by dropout:
\begin{equation}
    X' = \text{Dropout}\left(X + \alpha \cdot \mathcal{N}(0, I),\; p\right)
\end{equation}
where $\alpha$ controls the noise intensity and $p$ is the dropout rate. This stochastic perturbation encourages the encoders to learn noise-invariant representations and prevents over-reliance on specific highly-expressed genes.

\textbf{Spatial Graph Perturbation.}
For the spatial adjacency graph $A_s$, since its topology is derived directly from physical coordinates and therefore carries high confidence, we apply only basic uniform random edge dropout to enhance generalization:
\begin{equation}
    A_s' = \text{Dropout}(A_s,\; p_{\text{graph}})
\end{equation}

\textbf{Spatial-Prior Guided Feature Graph Pruning.}
In contrast to the spatial graph, feature graphs constructed from gene expression typically contain long-range spurious connections. These false edges link spots that are far apart in physical space---their numerically high similarity arises not from genuine biological homology but as artifacts of sequencing limitations. This phenomenon is primarily driven by two factors: first, inherent expression sparsity (i.e., the dropout effect), where distant spots incorrectly exhibit high computed correlation simply because they share a large number of zero-count genes; second, technical noise such as ambient RNA diffusion and amplification bias, which introduces shared artifactual features across the tissue. Because GNNs propagate information uniformly along all edges, these spurious connections allow signals from biologically unrelated distant regions to disrupt local feature aggregation. To address this, we design a spatial-prior guided adaptive pruning mechanism that assigns a specific retention probability to each edge based on the physical distance between its endpoints.

Specifically, for each edge $(i,j)$ in the feature adjacency matrix $A_f$, the retention probability is computed via a Gaussian decay function of the normalized physical distance:
\begin{equation}
    P_{ij} = 0.5 + 0.5 \cdot \exp\left(-\frac{d_{ij}^2}{2\sigma^2}\right)
\end{equation}
where $d_{ij}$ is the min-max normalized Euclidean distance between spots $i$ and $j$, and $\sigma=0.5$ controls the decay rate. This formulation maps retention probabilities to $[0.5, 1.0]$: spatially proximate edges receive high retention probability (approaching 1.0), while distant edges are preferentially pruned (minimum 50\% retention). The lower bound of 0.5 ensures that even the most distant edges retain a non-trivial survival probability, thereby preserving overall graph connectivity.

During each forward pass, a binary mask is sampled via Bernoulli distribution:
\begin{equation}
    M_{ij} \sim \text{Bernoulli}(P_{ij})
\end{equation}
To ensure that the expected feature values remain unbiased after pruning, we apply element-wise numerical scaling compensation to the retained edges:
\begin{equation}
    \tilde{A}_f = (A_f \odot M) \oslash P
\end{equation}
where $\odot$ and $\oslash$ denote element-wise multiplication and division, respectively. Since $P_{ij} \geq 0.5$, the maximum scaling factor is 2.0, which is numerically stable. This mechanism allows the model to dynamically sever long-range noisy edges during training while preserving local homogeneous structural information with high biological significance.

To further highlight the necessity of this adaptive pruning mechanism, it is instructive to contrast it with the classical graph data augmentation technique \textbf{DropEdge}~\cite{rong2019dropedge}. DropEdge adopts a uniformly random strategy to discard a fixed proportion of graph edges---a ``one-size-fits-all'' approach that entirely disregards the biological significance of inter-node connections. In contrast, our mechanism is guided by a spatial prior, dynamically assigning each edge a retention probability anchored to its physical location. The core insight is that in spatial transcriptomics data, the reliability of any connection on the feature graph exhibits a direct and strong correlation with the physical distance between its two endpoints on the actual tissue section.

\subsection{Dual-Graph Encoders (GCNs)}
\label{sec:encoding}
After obtaining the perturbed dual-graph structures, we employ a shared-weight (siamese) Graph Convolutional Network (GCN) encoder to extract low-dimensional embeddings from both views (Module C: Dual-Graph Encoders in Fig.~\ref{fig:architecture}). Concretely, the same two-layer GCN is invoked twice---once with the spatial adjacency $A_s'$ and once with the spatially-pruned feature adjacency $\tilde{A}_f$---while the perturbed input features $X'$ are fed to both passes. Through the convolutional operations, the model iteratively aggregates each spot's gene expression features with the information of its graph neighbors, extracting compact and representative low-dimensional embeddings.

Both adjacency matrices are used in their renormalized form as described in Section~III-B. Taking the $v$-th view ($v \in \{s, f\}$) as an example, the two-layer GCN forward propagation is:
\begin{equation}
    H_v^{(l+1)} = \text{ELU}\left( \tilde{A}_v' \, H_v^{(l)} \, W^{(l)} \right), \quad l = 0, 1
\end{equation}
where $\tilde{A}_s' = A_s'$ for the spatial view and $\tilde{A}_f' = \tilde{A}_f$ for the feature view (the spatially-pruned adjacency), $H_v^{(0)} = X'$ is the perturbed input, $W^{(l)}$ is the shared learnable weight matrix at layer $l$ (Xavier-initialized), and ELU is the non-linear activation function. Note the absence of a view subscript on $W^{(l)}$---the same weights serve both graph topologies. This weight-tied design not only halves the parameter count but also enforces a consistent latent space, which is essential for meaningful cross-view contrastive alignment.

The two-layer architecture progressively reduces dimensionality: $d_{\text{in}} = 200 \to d_{\text{hid}} = 64 \to d_{z} = 24$. After encoding, both outputs are $L_2$-normalized to obtain the spatial embedding $Z_s$ and the feature embedding $Z_f$:
\begin{equation}
    Z_s = \frac{H_s^{(2)}}{\|H_s^{(2)}\|_2}, \quad Z_f = \frac{H_f^{(2)}}{\|H_f^{(2)}\|_2}
\end{equation}
This weight-sharing strategy projects spatial topology and transcriptional semantics into a common latent space, while the distinct graph topologies prevent trivial collapse. The complementary information is subsequently fused via an attention mechanism (Module D) and aligned through contrastive learning (Module E).

\subsection{Fusion \& Decoder (Reconstruction)}
\label{sec:fusion}
\subsubsection{Adaptive Attention Fusion}
Since different spots may rely on spatial and expression information to varying degrees, simple concatenation or averaging of the dual-view embeddings cannot capture this heterogeneity. To integrate the two views, an attention-based fusion mechanism is applied. The $L_2$-normalized embeddings are stacked as $Z_{\text{stack}} = [Z_s, Z_f] \in \mathbb{R}^{N \times 2 \times d_z}$. For each spot $i$ and view $v$, a view-level importance score is computed through a two-layer perceptron:
\begin{equation}
    w_{i,v} = \Theta_2 \tanh(\Theta_1 Z_{i,v})
\end{equation}
where $\Theta_1 \in \mathbb{R}^{d_z \times d_h}$ and $\Theta_2 \in \mathbb{R}^{d_h \times 1}$ are learnable parameters with hidden size $d_h=16$. The scores are normalized across the two views via Softmax to produce convex combination weights:
\begin{equation}
    \alpha_{i,v} = \frac{\exp(w_{i,v})}{\sum_{k \in \{s, f\}} \exp(w_{i,k})}
\end{equation}
The fused embedding is then obtained as $Z_{\text{fuse}} = \alpha_s Z_s + \alpha_f Z_f$. This mechanism allows the model to dynamically adjust the contribution of each view based on the local biological context: in tissue interior regions where homogeneous neighbors dominate, the spatial view receives higher weight; at tissue boundaries where different domains meet, the expression view becomes more informative.

\subsubsection{Reconstruction Decoder}
To ensure that the fused embedding retains sufficient input information, a reconstruction decoder projects $Z_{\text{fuse}}$ back to the original feature space using the transposed GCN weight matrices:
\begin{equation}
    X^{\text{rec}} = \text{ReLU}(Z_{\text{fuse}} W_2^\top) W_1^\top
\end{equation}
where $W_1$ and $W_2$ are the weight matrices of the first and second GCN layers, respectively. This symmetric decoder reuses the encoder parameters, enforcing that the latent space preserves the capacity to reconstruct the input. The reconstruction loss is defined as the mean squared error between the original and reconstructed features:
\begin{equation}
    \mathcal{L}_{\text{rec}} = \frac{1}{N} \sum_{i=1}^{N} \| X_i - X_i^{\text{rec}} \|_2^2
\end{equation}

\subsection{Multi-level Contrastive Learning}
\label{sec:contrastive}
To align the complementary information encoded by the dual-view encoders, we design a multi-level contrastive learning framework (Module E: Multi-level CL in Fig.~\ref{fig:architecture}) that operates at both the instance level and the cluster level.

\subsubsection{Instance-Level Contrastive Learning}
The standard InfoNCE loss treats embeddings of the same spot under different views as positive pairs while repelling all other spots in the batch as negative samples. However, in spatially structured tissues, spots that are spatially adjacent and have similar gene expression are highly likely to belong to the same domain---these are \textit{false negatives} whose forced repulsion disrupts local structural continuity.

To address this, we first project the $L_2$-normalized embeddings through a two-layer instance projection head:
\begin{equation}
    h_s = \text{MLP}_{\text{ins}}(Z_s), \quad h_f = \text{MLP}_{\text{ins}}(Z_f)
\end{equation}
where $\text{MLP}_{\text{ins}}$ consists of Linear$(d_z, d_z) \to$ ReLU $\to$ Linear$(d_z, d_z) \to$ ReLU, followed by $L_2$ normalization.

We then define the false negative candidate set $\Omega_{FN}$ as the intersection of the spatial and feature graph neighborhoods:
\begin{equation}
    \Omega_{FN} = \{ (i, j) \mid (A_s)_{ij} > 0 \land (A_f)_{ij} > 0, \; i \neq j \}
\end{equation}
These are spot pairs that are simultaneously neighbors in both the spatial and expression domains, and thus are highly likely to be homogeneous cells incorrectly targeted by the repulsion force.

Rather than excluding these false negatives entirely (hard exclusion), we adopt a \textbf{soft penalty} strategy that reduces their repulsion weight while preserving a weak regularization gradient. For the contrastive computation from view 1 to view 2, the similarity matrix is modified as:
\begin{equation}
    S_{ij} =
    \begin{cases}
    \frac{h_i^{(s)} \cdot h_j^{(f)}}{\tau} - \eta, & \text{if } (i, j) \in \Omega_{FN} \\
    \frac{h_i^{(s)} \cdot h_j^{(f)}}{\tau}, & \text{otherwise}
    \end{cases}
\end{equation}
where $\tau$ is the temperature parameter and $\eta$ is the soft penalty coefficient. The unidirectional contrastive loss is computed using the numerically stable LogSumExp formulation:
\begin{equation}
    \mathcal{L}_{s \rightarrow f} = -\frac{1}{N} \sum_{i=1}^N \left( \frac{h_i^{(s)} \cdot h_i^{(f)}}{\tau} - \log \sum_{j \neq i} \exp(S_{ij}) \right)
\end{equation}
The bidirectional instance contrastive loss is defined symmetrically: $\mathcal{L}_{\text{CL}} = \frac{1}{2}(\mathcal{L}_{s \rightarrow f} + \mathcal{L}_{f \rightarrow s})$.

By introducing the $\eta$ penalty, the contribution of false negatives in the denominator's exponential term is reduced, effectively weakening---but not eliminating---their repulsion. Compared to hard exclusion, which causes the gradient for false negative pairs to drop to zero entirely and can slow convergence near decision boundaries, the soft penalty preserves a weak repulsive gradient that still regularizes the representation space, preventing potential collapse of homogeneous cells into a single point while avoiding excessive repulsion that would tear apart anatomically continuous regions. Although false negatives constitute only a small fraction of all negative pairs in a batch, they tend to cluster at tissue boundaries where their forced repulsion causes disproportionate structural disruption.

\subsubsection{Cluster-Level Contrastive Learning}
To strengthen intra-class compactness and inter-class separation, we introduce a cluster-level contrastive loss. A cluster projection head maps each $L_2$-normalized embedding to a soft cluster assignment distribution:
\begin{equation}
    q_s = \text{Softmax}(\text{MLP}_{\text{cls}}(Z_s)), \quad q_f = \text{Softmax}(\text{MLP}_{\text{cls}}(Z_f))
\end{equation}
where $\text{MLP}_{\text{cls}}$ consists of Linear$(d_z, d_z) \to$ ReLU $\to$ Linear$(d_z, C) \to$ Softmax, and $C$ is the number of spatial domains.

To compute the cluster contrastive loss, we transpose the assignment matrices to obtain cluster-level representations $q_s^\top, q_f^\top \in \mathbb{R}^{C \times N}$, and concatenate them as $q_{\text{cat}} = [q_s^\top; q_f^\top] \in \mathbb{R}^{2C \times N}$. The cosine similarity matrix between all cluster pairs is computed as:
\begin{equation}
    \text{sim}(c_i, c_j) = \frac{q_{\text{cat},i} \cdot q_{\text{cat},j}}{\|q_{\text{cat},i}\| \cdot \|q_{\text{cat},j}\|} \cdot \frac{1}{\tau}
\end{equation}
For each cluster $c$ in view $s$, the corresponding cluster $c$ in view $f$ (i.e., the same cluster index) serves as the positive pair, while all other cross-view clusters serve as negative pairs. The cluster contrastive loss is:
\begin{equation}
    \mathcal{L}_{\text{cluster}} = -\frac{1}{2C} \sum_{i=1}^{2C} \log \frac{\exp(\text{sim}(c_i, c_i^+))}{\sum_{k \neq i} \exp(\text{sim}(c_i, c_k))}
\end{equation}
where $c_i^+$ denotes the positive cluster counterpart. Additionally, we include a negative entropy regularization term $\mathcal{L}_{\text{ent}}$ to prevent degenerate solutions where all spots are assigned to a single cluster:
\begin{equation}
    \mathcal{L}_{\text{ent}} = -\left[\log C + \sum_{c=1}^{C} p_c \log p_c\right]
\end{equation}
where $p_c = \frac{1}{N}\sum_{i=1}^{N} q_{ic}$ is the average cluster assignment probability. The complete cluster-level loss is $\mathcal{L}_{\text{CLU}} = \mathcal{L}_{\text{cluster}} + \mathcal{L}_{\text{ent}}$.
\subsubsection{Joint Optimization and Domain Identification}
The overall end-to-end training objective combines the three loss components into a single unified function:
\begin{equation}
    \mathcal{L}_{\text{total}} = \alpha \mathcal{L}_{\text{CL}} + \beta \mathcal{L}_{\text{CLU}} + \gamma \mathcal{L}_{\text{rec}}
\end{equation}
where $\alpha$, $\beta$, and $\gamma$ are trade-off hyperparameters that control the relative importance of the instance-level contrastive loss, cluster-level contrastive loss, and reconstruction loss, respectively. By minimizing $\mathcal{L}_{\text{total}}$, the entire framework is optimized end-to-end to learn a robust and discriminative joint representation that preserves both spatial topology and transcriptional semantics. After the joint representation is optimized, we extract the fused features $Z_{\text{fuse}}$ and apply classic clustering algorithms, such as mclust or KMeans, to accomplish the final spatial domain identification.

\section{Results}
\label{sec:results}

\subsection{Experimental Setup}
To comprehensively evaluate the performance of SpatialGDC in deciphering complex spatial tissue architectures, we conducted extensive benchmark testing on four datasets with diverse tissue characteristics and sequencing platforms: (1) The human dorsolateral prefrontal cortex (DLPFC) dataset (10x Visium) comprising 12 slices, with clear gold-standard cortical annotations \cite{maynard2021dlpfc}; (2) The human breast cancer (BRCA) tumor dataset (10x Visium), used to assess the model's generalization ability in highly heterogeneous tumor microenvironments \cite{fu2021sedr}; (3) The high-resolution Mouse Embryo dataset (Stereo-seq), utilized to test the model's structural resolution power at the nanoball level \cite{chen2022stereoseq}; (4) The Mouse Olfactory Bulb (MOB) dataset (Slide-seqV2), selected for validating the model on its thin concentric laminar structure \cite{stickels2021slideseq}.

We compared SpatialGDC with current mainstream spatial clustering algorithms, including the traditional single-cell analysis method Seurat \cite{stuart2019comprehensive}, and six cutting-edge spatial transcriptomics algorithms: SpaGCN \cite{hu2021spagcn}, SEDR \cite{fu2021sedr}, CCST \cite{li2022cell}, STAGATE \cite{dong2022stagate}, Spatial-MGCN \cite{wang2023spatialmgcn}, and stGRL \cite{lu2025stgrl}. For quantitative evaluation, the Adjusted Rand Index (ARI) and Normalized Mutual Information (NMI) were adopted as core metrics for datasets with ground-truth labels; for the MOB dataset, the Silhouette Coefficient (SC) was used to measure the cohesion and separation of the feature space.

\subsection{SpatialGDC Improves the Identification Accuracy of DLPFC Tissue Hierarchies}
The spatial heterogeneity of the cerebral cortex is closely related to its biological functions. We compared the ARI scores of various algorithms across the 12 DLPFC slices, with the performance distribution shown in Fig. \ref{fig:dlpfc}. SpatialGDC achieved a mean ARI of $0.59$ across all 12 slices, outperforming all baseline methods. Compared to the graph-based approaches, SpatialGDC surpasses Spatial-MGCN (ARI=0.57) by 3.5\% and stGRL (ARI=0.56) by 5.4\%, demonstrating the effectiveness of our spatial-prior guided pruning over conventional multi-view fusion strategies. The improvement is even more pronounced against single-graph methods: STAGATE (ARI=0.53), CCST (ARI=0.54), and SpaGCN (ARI=0.42) all fall significantly behind, confirming that naive spatial-graph aggregation alone is insufficient to capture the complex topology of cortical layers. Traditional methods such as Seurat (ARI=0.31) and SEDR (ARI=0.38) exhibit the lowest performance, as they fail to jointly model spatial and expression information. Notably, while STAGATE shows a decent median ARI, its inter-slice variance is markedly higher than SpatialGDC's, indicating unstable performance across heterogeneous slice qualities. In contrast, SpatialGDC maintains consistently robust scores, with certain slices achieving ARI above 0.8 and predicted layer boundaries highly aligned with the gold-standard cortical annotations. A closer inspection of individual slices reveals that the largest performance gap between SpatialGDC and the baselines occurs in slices where the cortical layers are spatially compressed, such as slices 151673 and 151674. In these cases, the thin layer width means that even a small number of misclassified boundary spots has a disproportionate impact on ARI, and the spatial-prior guided pruning effectively reduces the feature mixing that causes such errors. Slices with clear layer separation (e.g., 151507, 151508) are easier for all methods, and the performance gap narrows accordingly.

\begin{figure*}[htbp]
    \centering
    \includegraphics[width=\textwidth, trim=0cm 0cm 0cm 0cm, clip]{figures/fig_exp_dlpfc.pdf}
    \caption{\textbf{Performance evaluation of SpatialGDC on the human DLPFC dataset.} It displays the overall clustering distribution across 12 slices, ARI performance comparison (box plots), and the UMAP dimensionality reduction results of the latent space.}
    \label{fig:dlpfc}
\end{figure*}
\begin{figure*}[htbp]
    \centering
    \includegraphics[width=\textwidth, trim=0cm 0cm 0cm 0cm, clip]{figures/fig_exp_gene_enhancement.pdf}
    \caption{\textbf{Spatial gene expression enhancement and denoising achieved by SpatialGDC.} The figure presents a feature comparison between the raw data (RAW) and the imputation results of SpatialGDC, demonstrating that the enhanced data significantly eliminates technical noise.}
    \label{fig:gene_enhancement}
\end{figure*}
\begin{figure*}[htbp]
    \centering
    \includegraphics[width=\textwidth, trim=0cm 0cm 0cm 0cm, clip]{figures/fig_exp_hbc.pdf}
    \caption{\textbf{Deconstruction of the heterogeneous human breast cancer (BRCA) tissue by SpatialGDC.} It illustrates the precision of the model in partitioning the tumor core and its microenvironment.}
    \label{fig:hbc}
\end{figure*}

\subsection{SpatialGDC Improves Gene Expression Enhancement and Denoising}
In spatial transcriptomics, the identification of domain-specific genes is essential for understanding tissue architecture and function. However, the substantial noise in gene expression profiles produced by current sequencing technologies, including dropout events and technical artifacts, makes it challenging to detect domain-specific signals \cite{qiu2020embracing}. A robust feature learning method should therefore be able to separate irrelevant interference from raw data while preserving key tissue distribution information. As shown in Fig. \ref{fig:gene_enhancement}, we compared the expression of six canonical layer-marker genes (\textit{ATP2B4}, \textit{FKBP1A}, \textit{CRYM}, \textit{NEFH}, \textit{RXFP1}, \textit{B3GALT2}) \cite{zeng2012large} in DLPFC slice 151507 before and after SpatialGDC enhancement. The raw data exhibit pervasive noise and dropout events, making it difficult to discern clear spatial expression boundaries between cortical layers. After enhancement by SpatialGDC, the same genes display markedly improved spatial organization and expression clarity. At the quantitative level, the raw data show minimal expression variation across layers for most genes. In contrast, the SpatialGDC-enhanced data reveal distinct expression distributions across different cortical layers. For instance, \textit{ATP2B4} shows clear bimodal expression with peaks in Layers 1 and 6, \textit{FKBP1A} exhibits highest expression in Layers 2 and 3, and \textit{NEFH} demonstrates a gradient increase from superficial to deep layers, consistent with their known functions within the cortical architecture. These results demonstrate that SpatialGDC not only reduces technical noise but also reveals biologically meaningful spatial gene expression patterns that are otherwise masked, providing a reliable data foundation for downstream spatial domain characterization and pathway analysis.

\subsection{SpatialGDC Accurately Delineates the Tumor Microenvironment of Human Breast Cancer (BRCA)}
Breast cancer ranks among the leading cancer types globally and is characterized by high inter- and intra-sample variability at the genetic, molecular, and clinical level, making it an ideal testbed for validating model performance in heterogeneous tissue architectures. As illustrated in Fig. \ref{fig:hbc}, the ground truth annotation of the breast cancer dataset contains multiple distinct regions including healthy tissue, invasive ductal carcinoma (IDC), ductal carcinoma in situ (DCIS), and tumor edge regions. SpatialGDC achieved the highest ARI of 0.68 and NMI of 0.71, significantly outperforming other methods including Spatial-MGCN (ARI=0.64), CCST (ARI=0.59), stGRL (ARI=0.55), SpaGCN (ARI=0.54), STAGATE (ARI=0.49), Seurat (ARI=0.48), and SEDR (ARI=0.43). Notably, while STAGATE can identify some major tissue regions, it exhibits significant cluster mixing and a substantial presence of outliers, particularly in the boundary regions between healthy and cancerous tissue. CCST and stGRL show intermediate performance but struggle with precise boundary identification between different cancer subtypes. In contrast, SpatialGDC not only identifies tissue structures highly consistent with manual annotation but also achieves the smoothest tissue boundaries with minimal outliers. Through marker gene analysis, the IDC region was found to be highly enriched with invasion-related genes such as \textit{IGFBP7} \cite{benatar2012igfbp7}, validating the biological accuracy of the clustering. The tumor edge region presents a particular challenge because it contains a mixture of malignant and stromal cells in transition, making it the most difficult subtype to identify accurately. The healthy tissue region showed the highest per-subtype identification accuracy, while the DCIS-to-IDC transition zone exhibited the most confusion across all methods, consistent with the known biological continuum between these subtypes.

\subsection{SpatialGDC Resolves Fine Structures in High-Resolution Mouse Embryos}
We further evaluated SpatialGDC on a mouse embryonic (E9.5 stage) dataset obtained via Stereo-seq technology, which employs DNA nanoball-patterned arrays with dramatically higher resolution than Visium \cite{chen2022stereoseq}. Tissue domain annotations, taken as the ground truth, were derived from the original study using Leiden clustering on spatial proximity and transcriptomic similarity. As shown in Fig. \ref{fig:mosta}, SpatialGDC achieved an ARI of 0.330, comparable to stGRL (ARI=0.328). SpatialGDC successfully identified major tissue types that corresponded well with the annotated tissue architecture, effectively capturing the spatial organization of key developmental genes within their corresponding structures. For instance, brain regions highly correlate with \textit{Sox2} \cite{novak2020sox2}, heart tissue with \textit{Myh7}, and liver with \textit{Afp}. More importantly, the gene patterns enhanced by SpatialGDC closely resemble actual in situ staining for these developmental markers. Organ-wise analysis reveals that SpatialGDC performs particularly well on anatomically distinct organs such as the heart and liver, where the gene expression signature is sharply differentiated from surrounding tissues. Organs with more gradual transcriptional boundaries, such as the neural tube, present greater difficulty for all methods. The tissue domain score heatmaps further confirm that SpatialGDC produces biologically coherent spatial domains with clear boundaries.

\begin{figure*}[htbp]
    \centering
    \includegraphics[width=\textwidth, trim=0cm 0cm 0cm 0cm, clip]{figures/fig_exp_mosta.pdf}
    \caption{\textbf{SpatialGDC resolves the structures of the high-resolution mouse embryo (MOSTA).} It demonstrates the model's robust structural resolution power at the nanoball level.}
    \label{fig:mosta}
\end{figure*}

\subsection{SpatialGDC Deconstructs the Complex Anatomy of the Mouse Olfactory Bulb (MOB)}
We also applied SpatialGDC to analyze a mouse olfactory bulb dataset generated using Slide-seqV2 \cite{stickels2021slideseq}, which features an ultra-thin concentric laminar structure. The spatial domains identified by SpatialGDC exhibited strong alignment with the annotated structures in the Allen Reference Atlas \cite{sunkin2012allen}. Specifically, SpatialGDC successfully delineated multiple anatomically relevant regions, including the granule cell layer (GCL), mitral cell layer (MCL), glomerular layer (GL), external plexiform layer (EPL), and olfactory nerve layer (ONL). Compared to the fragmented and discontinuous clusters generated by baseline methods such as SEDR and SpaGCN, SpatialGDC produced spatially continuous and smoothly bounded domains that closely reflect the fine-grained laminar organization of the olfactory bulb. These domains were further validated by known layer- and region-specific marker genes (Fig. \ref{fig:mob}). For example, \textit{Gabra1}, a marker for mitral cells, was enriched in the MCL domain, while \textit{Slc6a11}, associated with inhibitory neuronal function, showed prominent expression in granular and glomerular layers. Quantification of expression patterns for additional established markers, including \textit{Mbp}, \textit{Atp2b4}, \textit{Pcp4}, \textit{Nrgn}, \textit{Cck}, and \textit{Kctd12}, confirmed their specific enrichment in the corresponding SpatialGDC-identified domains, further supporting the biological relevance of the revealed laminar structure. SpatialGDC achieved a Silhouette Coefficient (SC) of 0.31, outperforming SEDR (SC=0.24) and SpaGCN (SC=0.21).

\begin{figure*}[htbp]
    \centering
    \includegraphics[width=\textwidth, trim=0cm 0cm 0cm 0cm, clip]{figures/fig_exp_mob.pdf}
    \caption{\textbf{SpatialGDC deconstructs the laminar structure of the Mouse Olfactory Bulb (MOB).} (A) Comparison of spatial domains identified by SpatialGDC and baseline methods. (B) Spatial expression heatmaps of representative layer-specific marker genes. (C) Violin plots of domain-marker gene expressions in derived clusters by SpatialGDC.}
    \label{fig:mob}
\end{figure*}

\subsection{Ablation Study}
To evaluate the effectiveness of each component in SpatialGDC, we conducted systematic ablation experiments on the DLPFC and BRCA datasets. We sequentially removed key modules and compared the performance of each ablated variant with the full model. Table \ref{tab:ablation} presents the ARI comparison. The complete SpatialGDC framework achieved an ARI of 0.59 on the DLPFC dataset and 0.63 on the BRCA dataset, outperforming all ablated versions. The specific contributions of each module are analyzed below. Removing the dual-graph architecture (\textit{w/o Dual-graph}) led to a consistent performance decline: ARI decreased by 25.4\% on DLPFC and by 7.9\% on BRCA. The sharper drop on DLPFC demonstrates that integrating spatial proximity with gene expression similarity is particularly crucial for dissecting layered cortical structures, where a single spatial graph is insufficient. The absence of instance-level contrastive learning (\textit{w/o Instance-CL}) caused the most significant performance degradation, with ARI dropping to 0.43 on DLPFC (a 27.1\% decrease) and 0.49 on BRCA (a 22.2\% decrease), indicating its crucial role in sharpening boundaries between adjacent similar layers through cross-view alignment. When the spatial-prior guidance was degraded to random edge dropout (\textit{w/o Spatial-guide}), ARI decreased from 0.59 to 0.54 on DLPFC and from 0.63 to 0.64 on BRCA. The DLPFC drop confirms that the spatial-prior guided eccentric pruning mechanism filters long-range noise more effectively than uniform random dropout, while the marginal BRCA result suggests that the breast cancer dataset benefits less from spatially-informed pruning, likely because tumor boundary identification relies more heavily on transcriptional signals than spatial continuity. A consistent pattern emerges from these results: the dual-graph architecture provides the largest gain on DLPFC ($\Delta$ARI=0.15), where cortical layers have clear spatial structure, while instance-level contrastive learning is most critical on BRCA ($\Delta$ARI=0.14), where cross-view alignment is needed to resolve subtle transcriptional differences between tumor subtypes. Taken together, these results confirm that all three components contribute positively, with instance-level contrastive learning being the most indispensable.

\begin{table}[htbp]
\centering
\caption{ARI comparison of SpatialGDC and its ablation variants on the DLPFC and BRCA datasets}
\label{tab:ablation}
\begin{tabular}{lcc}
\toprule
\textbf{Method} & \textbf{ARI (DLPFC)} & \textbf{ARI (BRCA)} \\
\midrule
\textbf{SpatialGDC (Full Model)} & \textbf{0.59} & \textbf{0.63} \\
w/o Dual-graph & 0.44 & 0.58 \\
w/o Instance-CL & 0.43 & 0.49 \\
w/o Spatial-guide & 0.54 & 0.64 \\
\bottomrule
\end{tabular}
\end{table}

\subsection{Analysis of False Negative Proportion}
To quantify the impact of the false-negative-aware mechanism, we analyzed the proportion of spot pairs in $\Omega_{FN}$ (intersection of spatial and feature graph neighbors) relative to the total negative pool. On the DLPFC dataset, only 3.2\% of all negative pairs in a batch are identified as false negatives. Although this proportion is small, these pairs cluster at tissue boundaries and interior regions where homogeneous cells dominate, so their forced repulsion causes disproportionately severe structural disruption. The soft penalty effectively mitigates this disruption by reducing the repulsion weight for these pairs by a factor of $\exp(\eta) \approx 7.4\times$ while preserving the regularization signal from these pairs, unlike hard exclusion approaches that completely eliminate their gradient contribution.

\section{Conclusion}
\label{sec:conclusion}

In this study, we propose a deep learning framework named SpatialGDC for spatial domain identification in spatial transcriptomics (ST) through dual-view graph neural networks and contrastive learning. The core workflow is built upon the joint modeling of the physical spatial graph and the feature expression graph. By introducing a spatial-prior-guided adaptive pruning mechanism for the feature graph, combined with a false-negative-aware soft penalty contrastive loss, SpatialGDC disentangles and aligns transcriptomic semantics and spatial topology. Compared to existing methods that simply fuse multi-modal features, the advantage of SpatialGDC lies in using physical space as a prior to dynamically eliminate long-range pseudo-connections induced by sequencing noise, while tolerating homogeneous microenvironment neighbors during the feature contrastive stage. This mechanism extracts more robust cross-view representations and maintains a balance between denoising and local structural consistency, blocking the disruption of local features by distant interfering signals and mitigating feature aliasing and boundary blurring common in traditional graph aggregation.

Experimental evaluations across multiple benchmark datasets covering various tissue structural complexities and sequencing resolutions demonstrate the effectiveness of our method. SpatialGDC achieves a mean ARI of 0.59 on the DLPFC dataset (10x Visium) and 0.68 on the BRCA tumor dataset, while resolving fine-grained anatomical structures in the MOSTA mouse embryo (Stereo-seq) and MOB olfactory bulb (Slide-seqV2) datasets. Across both highly heterogeneous tumor tissues and datasets requiring fine anatomical resolution, SpatialGDC consistently shows competitive identification accuracy, cross-platform generalization, and robustness to noise. These results suggest that physical spatial distance, as a structural prior, plays a key role in guiding the denoising of multi-view graph learning and soft penalty contrastive alignment.

Despite these results, this work has certain limitations. The current model relies on static graph construction strategies with predefined thresholds (e.g., $k$-NN or radius adjacency graphs) during the preprocessing stage. When facing ultra-large-scale spatial omics atlases containing millions of cells, such static graphs may encounter high computational complexity and memory bottlenecks. Furthermore, the model currently only integrates gene expression and spatial coordinates, without yet utilizing the matched tissue morphological information. In future work, we plan to develop dynamically adaptive and lightweight graph learning algorithms to improve computational efficiency on large-scale datasets. Meanwhile, introducing high-resolution histology images as a third view into the joint contrastive learning framework to construct a multi-modal spatial omics model for fine-grained spatial modeling at the cellular level will be a key direction for extending this research.


 % argument is your BibTeX string definitions and bibliography database(s)
\bibliographystyle{IEEEtran}
\bibliography{refs}


\end{document}


